% Section 6, part 2.
Throughout \cref{lem:mod-laplace,lem:mod-scale,prop:mod-Rtr}, $\kappa\neq 0$ is a rational number and $x_\kappa:=\kappa\Wv$; the curve $(\bbP^1,x_\kappa,y=p)$ has the branch points $\varsigma_j$ of \cref{def:mod-curve}. The cases used are $\kappa=4$ and $\kappa=1$. The letter $\vartheta$ denotes the Laplace variable of \cite[(2-31), (2-32)]{Dunin-Barkowski:2015caa} (called $\zeta$ there).

\begin{lem}[Formal Laplace series]
\label{lem:mod-laplace}
Let $c$ be a branch point of $x_\kappa$ and $w$ a local coordinate at $c$ with $x_\kappa=x_\kappa(c)-w^2/2$. For a $1$-form $\omega=\sum_{m\geq -M}a_mw^m\,\dd w$ meromorphic at $c$ with $\Res_c\omega=0$ put
\begin{equation}
\label{eq:mod-laplace}
\begin{gathered}
\mathsf L_c[\omega](\vartheta):=\sum_ma_m\,\mathsf G_m(\vartheta),\\
\mathsf G_{2n}:=(2n-1)!!\,\vartheta^{-n}\ (n\geq 0),\qquad \mathsf G_{-2n}:=\frac{(-1)^n\vartheta^n}{(2n-1)!!}\ (n\geq 1),\qquad \mathsf G_m:=0\ (m\text{ odd}),
\end{gathered}
\end{equation}
a formal Laurent series in $\vartheta^{-1}$ (the coefficient $a_{-1}$ vanishes). Then:
\begin{enumerate}
\item if $F$ is meromorphic at $c$ and $F\,\dd x_\kappa$ has zero residue at $c$, then $\mathsf L_c[\dd F]=-\vartheta\,\mathsf L_c[F\,\dd x_\kappa]$;
\item $w$ is determined by $x_\kappa$ up to $w\mapsto -w$, which multiplies $\mathsf L_c$ by $-1$;
\item let $k\in\bbC^\times$ and $v=k^5v'$, so that $x_{\kappa,v}(kp')=k^4x_{\kappa,v'}(p')$ and $\varsigma_j(v)=k\varsigma_j(v')$, and normalise the coordinates $w_j$ at $\varsigma_j$ by $\frac{\dd w_j}{\dd p}(\varsigma_j)=\sqrt{-5\kappa}\,\varsigma_j$ with one square root for all $v$; then $w_j(v)(kp')=k^2w_j(v')(p')$ and $\mathsf L_{\varsigma_j(v)}[p^m\dd p](\vartheta)=k^{m-1}\,\mathsf L_{\varsigma_j(v')}[p'^m\dd p'](k^4\vartheta)$;
\item if $\omega$ is holomorphic at $c$, then $\vartheta\frac{\dd}{\dd\vartheta}\mathsf L_c[\omega]=\vartheta\,\mathsf L_c\big[(x_\kappa-x_\kappa(c))\,\omega\big]+\frac12\mathsf L_c[\omega]$.
\end{enumerate}
\end{lem}

The integrals $\int\omega\,e^{(x_\kappa-x_\kappa(c))\vartheta}$ in \cite[(2-31), (2-32)]{Dunin-Barkowski:2015caa} are read as $\sqrt{2\pi/\vartheta}\,\mathsf L_c[\omega]$. For the diagonal entries of (2-31) the integrand has the polar part $\dd w/w^2$, and $\mathsf G_{-2}=-\vartheta$ produces exactly the entry $1$ of $R^{-1}=1+O(\vartheta^{-1})$; this is the prescription of \cite[\S3.1, \S3.5]{DuninBarkowski:2012bw} (the singular part $\dd w\,\dd w'/(w-w')^2$ gives the Kronecker delta, the regular part is transformed). (For orientation only: $\sqrt{2\pi/\vartheta}\,\mathsf L_c[\omega]$ is the asymptotic expansion, as $\vartheta\to+\infty$, of $\int\omega\,e^{-\vartheta w^2/2}$ along the real $w$-line deformed near $0$, by Watson's lemma; this is not used.)

\begin{proof}
All four statements are $\bbC$-linear in $\omega$ (resp.\ $F$), so it suffices to prove them for monomials. (i) For $F=w^m$, $m\neq -2$: $\dd F=mw^{m-1}\dd w$ and $F\,\dd x_\kappa=-w^{m+1}\dd w$, so the claim is $m\mathsf G_{m-1}=\vartheta\mathsf G_{m+1}$. For $m=2n+1\geq 1$ both sides are $(2n+1)!!\,\vartheta^{-n}$; for $m=-(2n+1)\leq -1$ both are $(-1)^n\vartheta^{n+1}/(2n-1)!!$; for even $m$ both vanish. The monomial $F=w^{-2}$ is excluded, since $F\,\dd x_\kappa=-w^{-1}\dd w$ has a residue. (ii) The coordinate is fixed by $w^2=2(x_\kappa(c)-x_\kappa)$ up to sign, and $w^m\dd w\mapsto(-1)^{m+1}w^m\dd w$ under $w\mapsto -w$, while only even $m$ contribute. (iii) Both $w_j(v)(kp')$ and $k^2w_j(v')(p')$ are coordinates at $\varsigma_j(v')$ in the variable $p'$ whose squares equal $2(x_{\kappa,v}(\varsigma_j(v))-x_{\kappa,v}(kp'))=2k^4(x_{\kappa,v'}(\varsigma_j(v'))-x_{\kappa,v'}(p'))$; their derivatives at $p'=\varsigma_j(v')$ are $k\sqrt{-5\kappa}\,\varsigma_j(v)$ and $k^2\sqrt{-5\kappa}\,\varsigma_j(v')$, which agree, so the two coordinates agree. Write $p'^m\dd p'=\sum_nb_nw'^n\dd w'$ with $w'=w_j(v')$. Then $p^m\dd p=k^{m+1}p'^m\dd p'=\sum_nk^{m-1-2n}b_nw^n\dd w$ with $w=w_j(v)=k^2w'$, and $\mathsf G_{2s}(\vartheta)=k^{4s}\mathsf G_{2s}(k^4\vartheta)$, $\mathsf G_{-2s}(\vartheta)=k^{-4s}\mathsf G_{-2s}(k^4\vartheta)$ give the claim. (iv) On $w^{2n}\dd w$ the left side is $-n(2n-1)!!\,\vartheta^{-n}$; since $(x_\kappa-x_\kappa(c))w^{2n}\dd w=-\frac12w^{2n+2}\dd w$, the right side is $-\frac12(2n+1)!!\,\vartheta^{-n}+\frac12(2n-1)!!\,\vartheta^{-n}$, the same; odd powers give $0$ on both sides.
\end{proof}

\begin{lem}[Scaling of the curve]
\label{lem:mod-scale}
\begin{enumerate}
\item For constants $\kappa,\mu\neq 0$ let $\omega^{(\kappa,\mu)}_{h,n}$ and $F^{(\kappa,\mu)}_h$ be the correlators and free energies of $(\bbP^1,\kappa\Wv,\mu p,\omega_{0,2})$. Then $\omega^{(\kappa,\mu)}_{h,n}=(\kappa\mu)^{2-2h-n}\omega_{h,n}$ for $2h-2+n>0$, and $F^{(\kappa,\mu)}_h=(\kappa\mu)^{2-2h}\FTR_h$ for $h\geq 2$.
\item For $k\in\bbC^\times$ let $m_k(p):=kp$. Then $m_k^*\,\omega^{(k^5v)}_{h,n}=k^{5(2-2h-n)}\,\omega^{(v)}_{h,n}$ and $\FTR_h(k^5v)=k^{5(2-2h)}\FTR_h(v)$, where the superscript is the parameter of the curve. Consequently a pairing of $\omega_{h,n}$ with the legs $p^{m_1},\dots,p^{m_j}$ in $j$ slots and with $-\log p$ in the other $n-j$ slots is a function $C(v)$ with $C(k^5v)=k^{5e}C(v)$, $e=2-2h-n+\sum_lm_l/5$.
\end{enumerate}
\end{lem}

\begin{proof}
(i) The branch points and the involutions $\iota_b$ are those of $(\Wv,p)$, $\omega_{0,2}$ is unchanged, and the kernel becomes $K_b/(\kappa\mu)$. Suppose the claim holds for all $(h',n')$ with $2h'-2+n'<2h-2+n+1$, the exponent of $\omega_{0,2}$ being $0=2-0-2$. Every term of \eqref{eq:mod-kernel} for $\omega_{h,n+1}$ is one kernel times $\omega_{h-1,n+2}$, of exponent $2-2h-n$, or times a product $\omega_{h_1,|I_1|+1}\omega_{h_2,|I_2|+1}$, of exponent $(1-2h_1-|I_1|)+(1-2h_2-|I_2|)=2-2h-n$; the total exponent is $1-2h-n=2-2h-(n+1)$, which is the induction step. For the free energy, $\Phi_{y\dd x}$ becomes $\kappa\mu\,\Phi_{y\dd x}$, and $\kappa\mu(\kappa\mu)^{1-2h}=(\kappa\mu)^{2-2h}$. (ii) The map $m_k$ is an automorphism of $\bbP^1$ with $m_k^*W_{k^5v}=k^4\Wv$, $m_k^*p=kp$ and $m_k^*\omega_{0,2}=\omega_{0,2}$; it maps the branch points of $\Wv$ to those of $W_{k^5v}$. The recursion \eqref{eq:mod-kernel} is natural under isomorphisms of curves, so $m_k^*\omega^{(k^5v)}_{h,n}$ is the correlator of $(\bbP^1,k^4\Wv,kp,\omega_{0,2})$, which is $k^{5(2-2h-n)}\omega^{(v)}_{h,n}$ by (i); likewise for $\FTR_h$. For the pairings, pull the residues back along $m_k$; a leg $p_l^{m_l}$ contributes $k^{m_l}$, and a ray is mapped to a ray, along which the pairing does not depend on the choice (\cref{def:mod-pair}).
\end{proof}

\begin{prop}[The $R$-matrix of the curve]
\label{prop:mod-Rtr}
Let $v\neq 0$ and take at the branch points of $x_\kappa$ the coordinates $w_j$ of \cite[\S2.6]{Dunin-Barkowski:2015caa}, $x_\kappa=x_\kappa(\varsigma_j)-w_j^2/2$, normalised by $\frac{\dd w_j}{\dd p}(\varsigma_j)=\sqrt{-5\kappa}\,\varsigma_j$ with one square root for all $j$. Let $e^{\mathrm{TR}}_j$, $\Delta_j$ and $\xi^j$ be the corresponding data (2-30), (2-34) of \cite[\S2.6]{Dunin-Barkowski:2015caa}: $\Delta_j^{-1/2}=\frac{\dd y}{\dd w_j}(\varsigma_j)$ and $\xi^j(q)=\int_\infty^q\omega_{0,2}(w_j,\cdot)/\dd w_j|_{w_j=0}$; and let $R_{\mathrm{TR}}$ be the matrix series defined by (2-31) read through \cref{lem:mod-laplace}, $R_{\mathrm{TR}}^{-1}(\vartheta^{-1})^j{}_i:=-\vartheta^{-1}\mathsf L_{\varsigma_j}[\dd\xi^i]$ (upper index the row). Then, as formal series in $\vartheta^{-1}$:
\begin{enumerate}
\item $R_{\mathrm{TR}}^{-1}(\vartheta^{-1})^j{}_i=\Mo(t)_{ij}$ with $t=-1/(\kappa\varsigma^4\vartheta)$, for all $i,j$; consequently
\begin{equation}
\label{eq:mod-Rtr}
R_{\mathrm{TR}}(\psi)=\Mo\big(\psi/(\kappa\varsigma^4)\big)
\end{equation}
in the basis $e^{\mathrm{TR}}_j$, and $R_{\mathrm{TR}}$ is symplectic in this orthonormal basis;
\item (2-32) holds: $\sum_kR_{\mathrm{TR}}^{-1}(\vartheta^{-1})^i{}_k\Delta_k^{-1/2}=\mathsf L_{\varsigma_i}[\dd y]$, both sides being $\varsigma^{-1}\Rhat^{\mathrm m,(i)}_0(t)/\sqrt{-5\kappa}$. In the frozen basis $\bar e_j:=e^{\mathrm{TR}}_j/\sqrt{-5\kappa}$ the unit $\sum_k\Delta_k^{-1/2}e^{\mathrm{TR}}_k$ is $\hat 1:=\sum_k\bar e_k/\varsigma_k$, and the dilaton leaves of the graph formula carry the translation $\psi\big(\hat 1-\Mo^{-1}(\psi/(\kappa\varsigma^4))\hat 1\big)=\psi\,\hat 1-\psi R_{\mathrm{TR}}^{-1}(\psi)\hat 1$ of \cite[Def.~2.13]{PPZ15}.
\end{enumerate}
\end{prop}

The sign of the dilaton leaves is that of \cite[\S2.3]{DuninBarkowski:2012bw} and of \cite[Prop.~2.12, Def.~2.13]{PPZ15}; \cite[\S2.3]{Dunin-Barkowski:2015caa} prints the opposite sign. The three sources treat the dilaton leaves in different forms: in \cite[\S2.3]{DuninBarkowski:2012bw} they are marked points, and a vertex is an integral of classes $\psi$ over the moduli space with these points; \cite[\S2.3]{Dunin-Barkowski:2015caa} takes the classes $\psi$ of the legs and edges on that space and then forgets the dilaton points by push-forward; \cite[Defs.~2.2, 2.5]{PPZ15} forget them first and take the classes $\psi$ on the smaller space. By \cref{lem:set-pushins}~(ii), applied at every vertex, the last two forms give the same classes, whose integrals are the sums over the dilaton leaves of the vertex integrals of the first form, when the translation is the series of the weights of the dilaton leaves; so the weights of \cite[\S2.3]{DuninBarkowski:2012bw} fix the sign of the translation in the push-forward form (\cref{prop:mod-Rtr}~(ii)). With legs $R^{-1}(\psi)$ the sign is not a convention: it is the one for which the action preserves the unit (\cref{sec:sgn-R}).

% Section 6, part 3.
\begin{proof}[Proof of \cref{prop:mod-Rtr}]
Write $\sqrt{-5\kappa}$ for the root of the statement and $\mathsf L_j:=\mathsf L_{\varsigma_j}$; indices run over $\bbZ/5$.

\emph{Step 1 (integration by parts).} By (2-34), $\dd\xi^i(q)=\frac{\dd p}{\dd w_i}(\varsigma_i)\,\dd q/(\varsigma_i-q)^2$, so $\xi^i(q)=-1/\big(\sqrt{-5\kappa}\,\varsigma_i(q-\varsigma_i)\big)$. This function has at most a simple pole at $\varsigma_j$, so $\xi^i\dd x_\kappa$ is holomorphic there, and \cref{lem:mod-laplace}~(i) gives $R_{\mathrm{TR}}^{-1}(\vartheta^{-1})^j{}_i=-\vartheta^{-1}\mathsf L_j[\dd\xi^i]=\mathsf L_j[\xi^i\dd x_\kappa]$.

\emph{Step 2 (partial fractions).} Since $\dd x_\kappa=\kappa(p^5-v)\,\dd p/p^2$ and $(p^5-v)/(p-\varsigma_i)=\sum_{k=0}^4\varsigma_i^kp^{4-k}$, we have $\xi^i\dd x_\kappa=-\frac{\kappa}{\sqrt{-5\kappa}\varsigma_i}\sum_{k=0}^4\varsigma_i^kp^{2-k}\dd p$. With the local letters $J^{(j)}_m:=\mathsf L_j[p^m\dd p]$, $\tilde J^{(j)}_m:=\sqrt{-5\kappa}\,J^{(j)}_m$ and $-\kappa/(-5\kappa)=1/5$:
\begin{equation}
\label{eq:mod-Rtr-pf}
R_{\mathrm{TR}}^{-1}(\vartheta^{-1})^j{}_i=\frac15\sum_{k=0}^4\varsigma_i^{k-1}\,\tilde J^{(j)}_{2-k}.
\end{equation}

\emph{Step 3 (the five series are model rows).} Fix $j$ and write $J_m:=J^{(j)}_m$. Let $z:=-1/(\kappa\vartheta)$, a formal variable with $\vartheta x_\kappa=-\Wv/z$, and $t:=z/\varsigma^4=-1/(\kappa\varsigma^4\vartheta)$. \emph{(Mod2).} \Cref{lem:mod-laplace}~(i) with $F=p^m$, which is holomorphic at $\varsigma_j\neq 0$, gives $mJ_{m-1}=-\vartheta\kappa(J_{m+3}-vJ_{m-2})$, that is $mzJ_{m-1}-J_{m+3}+vJ_{m-2}=0$; in particular $J_3=vJ_{-2}$. \emph{(Mod1).} By \cref{lem:mod-laplace}~(iii) with $k=\varsigma$, $v'=1$, we have $J_m(v;\vartheta)=\varsigma^{m-1}J_m(1;\varsigma^4\vartheta)$, so the coefficient of $\vartheta^{-n}$ in $J_m$ is a constant times $\varsigma^{m-1-4n}$; differentiating termwise, $5v\,\partial_vJ_m=(m-1)J_m+4\vartheta\partial_\vartheta J_m$. By \cref{lem:mod-laplace}~(iv), with $x_\kappa-x_\kappa(\varsigma_j)=\kappa(p^4/4+v/p)-\frac54\kappa\varsigma_j^4$,
\[
\vartheta\partial_\vartheta J_m=\vartheta\kappa\big(\tfrac14J_{m+4}+vJ_{m-1}\big)-\tfrac54\vartheta\kappa\varsigma_j^4J_m+\tfrac12J_m .
\]
Eliminating $J_{m+4}=(m+1)zJ_m+vJ_{m-1}$ by (Mod2) and using $\kappa\vartheta=-1/z$ gives $\vartheta\partial_\vartheta J_m=-\frac{m+1}4J_m-\frac{5v}{4z}J_{m-1}+\frac{5\varsigma_j^4}{4z}J_m+\frac12J_m$, hence $5v\,\partial_vJ_m=-\frac{5v}{z}J_{m-1}+\frac{5\varsigma_j^4}{z}J_m$, that is $z\partial_vJ_m=-J_{m-1}+J_m/\varsigma_j$ (as $\varsigma_j^5=v$). Put $\tilde Y_m:=e^{-\Wv(\varsigma_j)/z}J_m$, a formal symbol on which $\partial_v$ acts by the Leibniz rule. Since $\partial_v[\Wv(\varsigma_j(v))]=1/\varsigma_j$ (because $\Wv'(\varsigma_j)=0$ and $\partial_v\Wv=1/p$), we get $z\partial_v\tilde Y_m=-\tilde Y_{m-1}$. The five series $Y_b:=\varepsilon_b\tilde Y_{\nu_b}$, $b\in\bbZ/5$, therefore satisfy $z\partial_vY_b=b^{\mathrm m}_bY_{b+1}$: indeed $z\partial_vY_0=-\tilde Y_{-1}=Y_1$, $z\partial_vY_1=\tilde Y_{-2}=-Y_2/v$ (as $Y_2=-\tilde Y_3=-v\tilde Y_{-2}$ by (Mod2)), $z\partial_vY_2=\tilde Y_2=Y_3$, $z\partial_vY_3=-\tilde Y_1=-Y_4$ and $z\partial_vY_4=-\tilde Y_0=-Y_0$. This is the derivation of \cref{prop:set-modrows}~(a) from (Mod1) and (Mod2), now for these formal series. Write $Y_b=e^{U_j/z}\varsigma^{\nu_b-1}F_b$ with $U_j:=-\Wv(\varsigma_j)=-\frac54\zeta^j\varsigma^4$ and $F_b:=\varepsilon_b\varsigma^{1-\nu_b}J_{\nu_b}$. By the homogeneity $J_m(v;\vartheta)=\varsigma^{m-1}J_m(1;\varsigma^4\vartheta)$ used for (Mod1), $F_b=\varepsilon_bJ_{\nu_b}(1;\varsigma^4\vartheta)$ depends only on $\varsigma^4\vartheta=-1/(\kappa t)$, so $F_b$ is a power series in $t$, and the computation of \cref{prop:set-modrows}~(b), $z\partial_v\big[e^{U_j/z}\varsigma^{\nu-1}F(t)\big]=e^{U_j/z}\varsigma^{\nu-2}\mathcal T^{(j)}_{\nu-1}F$, shows that $(F_b)_b$ solves the $T$-system of colour $j$. \emph{Initial values.} $J_m=\varsigma_j^m\frac{\dd p}{\dd w_j}(\varsigma_j)+O(\vartheta^{-1})=\varsigma_j^{m-1}/\sqrt{-5\kappa}+O(\vartheta^{-1})$, and $(\varsigma_j/\varsigma)^{\nu_b-1}=\zeta^{-j(\nu_b-1)}=\zeta^{j(b+1)}$ since $\nu_b-1\equiv -(b+1)\pmod 5$; so $\sqrt{-5\kappa}F_b(0)=\varepsilon_b\zeta^{j(b+1)}=\Rhat^{\mathrm m,(j)}_b(0)$ (\cref{def:set-M0}). \emph{The model column of colour $j$.} With $s:=\zeta^{-j}t$, so that $t\,\dd/\dd t=s\,\dd/\dd s$, one has $\mathcal T^{(j)}_k\big[\zeta^{j(b+1)}\Fs_b(\zeta^{-j}t)\big]=\zeta^{j(b+2)}\big(\mathcal T^{(0)}_k\Fs_b\big)(s)$; so the colour-$0$ relations of \cref{prop:set-modrows}~(b) for $\Fs$ become the colour-$j$ relations for $\Rhat^{\mathrm m,(j)}_b(t)=\zeta^{j(b+1)}\Fs_b(\zeta^{-j}t)$, the cycle $4\to 0$ being closed by $\zeta^{5j}=1$. \emph{Uniqueness.} By the uniqueness clause of \cref{prop:set-modrows}~(b) (pivot $4n\zeta^{4j}$), a power-series solution of the colour-$j$ system is determined by the constant term of its component $b=0$. Hence $\sqrt{-5\kappa}F_b=\Rhat^{\mathrm m,(j)}_b(t)$ for all $b$, that is $\tilde J^{(j)}_{m_b}=\varepsilon_b\varsigma^{m_b-1}\Rhat^{\mathrm m,(j)}_b(t)$ for $(m_0,\dots,m_4)=(0,-1,-2,2,1)$, $m_b\equiv -b\pmod 5$ (for $b=2$ use $\tilde J_{-2}=\tilde J_3/v$).

\emph{Step 4 (assembly).} As $b$ runs over $\bbZ/5$, $k:=2-m_b$ runs over $0,\dots,4$, and $\varsigma_i^{k-1}\varsigma^{m_b-1}=(\varsigma_i/\varsigma)^{1-m_b}=\zeta^{-i(b+1)}$. Inserting Step 3 into \eqref{eq:mod-Rtr-pf},
\[
R_{\mathrm{TR}}^{-1}(\vartheta^{-1})^j{}_i=\frac15\sum_b\varepsilon_b\zeta^{-i(b+1)}\,\Rhat^{\mathrm m,(j)}_b(t)=\sum_b(\Psi_0^{-1})_{ib}\Rhat^{\mathrm m,(j)}_b(t)=\Mo(t)_{ij},
\]
since $(\Psi_0)_{bi}=\varepsilon_b\zeta^{i(b+1)}$ has the inverse $(\Psi_0^{-1})_{ib}=\frac15\varepsilon_b\zeta^{-i(b+1)}$ ($\varepsilon_b^2=1$ and $\sum_b\zeta^{(j-i)(b+1)}=5\delta_{ij}$).

\emph{Step 5 (inverse).} As matrices with the upper index as row, Step 4 reads $R_{\mathrm{TR}}^{-1}(\psi)=\Mo\big(-\psi/(\kappa\varsigma^4)\big)^T$ with $\psi=\vartheta^{-1}$. By \cref{lem:set-M0sympl}, $\Mo(-t)^T=\Mo(t)^{-1}$, so $R_{\mathrm{TR}}(\psi)=\Mo(\psi/(\kappa\varsigma^4))$, and it is symplectic by the same lemma. This proves (i).

(ii) By (2-30), $\Delta_k^{-1/2}=\frac{\dd p}{\dd w_k}(\varsigma_k)=1/(\sqrt{-5\kappa}\varsigma_k)$. By (i), using $1/\varsigma_k=\zeta^k/\varsigma$ and $\sum_k\zeta^{-kb}=5\delta_{b,0}$, the left side of (2-32) is
\[
\sum_k\frac{\Mo(t)_{ki}}{\sqrt{-5\kappa}\,\varsigma_k}=\frac{1}{\sqrt{-5\kappa}\,\varsigma}\sum_b\varepsilon_b\Big[\frac15\sum_k\zeta^{-kb}\Big]\Rhat^{\mathrm m,(i)}_b(t)=\frac{\Rhat^{\mathrm m,(i)}_0(t)}{\sqrt{-5\kappa}\,\varsigma},
\]
and the right side is $\mathsf L_i[\dd p]=J^{(i)}_0=\tilde J^{(i)}_0/\sqrt{-5\kappa}=\Rhat^{\mathrm m,(i)}_0(t)/(\sqrt{-5\kappa}\varsigma)$ by Step 3. So (2-32) holds to all orders. The unit is $\sum_k\Delta_k^{-1/2}e^{\mathrm{TR}}_k=\sum_k\bar e_k/\varsigma_k$. Given (2-32), the graph formula of \cite[Thm.~4.1]{DuninBarkowski:2012bw}, as recalled in \cite[\S2.6]{Dunin-Barkowski:2015caa}, decorates the dilaton leaves by the translation $\psi\hat 1-\psi R_{\mathrm{TR}}^{-1}(\psi)\hat 1$ of \cite[Def.~2.13]{PPZ15}: in \cite[\S2.3]{DuninBarkowski:2012bw} a dilaton leaf is a marked point of the moduli space of its vertex carrying $\psi^k$, $k\geq 2$, with the weight $\coef{z^{k-1}}\bigl(-R(-z)\bigr)$ applied to the unit, which summed over $k$ is $\psi\bigl(1-R(-\psi)\bigr)$ applied to the unit; by \cref{lem:set-pushins}~(ii), applied at every vertex, these marked points give the same vertex integrals as the translation of \cite[Def.~2.5]{PPZ15} by this series, whose points are forgotten by push-forward; and the matrix $R(-z)$ of \cite[\S2.3]{DuninBarkowski:2012bw}, which also decorates the ordinary leaves there, is the matrix $R^{-1}(z)$ of \cite[\S2.3]{Dunin-Barkowski:2015caa}. Finally $R_{\mathrm{TR}}^{-1}(\psi)=\Mo(\psi/(\kappa\varsigma^4))^{-1}$ by (i).
\end{proof}

\begin{proof}[Proof of \cref{thm:mod-tr}]
\emph{Step 1 (the quoted theorem).} \cite[Thm.~4.1]{DuninBarkowski:2012bw}, as recalled in \cite[\S2.6, (2-30)--(2-35)]{Dunin-Barkowski:2015caa}, with the action of \cite[\S2.3]{Dunin-Barkowski:2015caa} (legs $R^{-1}(\psi)$ and edges as in \cite[Def.~2.2]{PPZ15}, dilaton leaves with the translation of \cite[Def.~2.13]{PPZ15}; by \cref{lem:set-pushins}~(ii) its classes are those of the push-forward form of \cite[Def.~2.5]{PPZ15}, see the paragraph after \cref{prop:mod-Rtr}), gives the following for a compact curve whose branch points are simple with $\dd y\neq 0$ there, with the Bergman kernel: if the matrix $R$ defined by (2-31) is symplectic in the orthonormal basis and (2-32) holds, then for $2h-2+n>0$
\begin{equation}
\label{eq:mod-235}
\omega_{h,n}=\sum_{i_1,\dots,i_n}\sum_{d_1,\dots,d_n\geq 0}\int_{\overline{\mathcal M}_{h,n}}\alpha_{h,n}(e_{i_1},\dots,e_{i_n})\prod_{l=1}^n\psi_l^{d_l}\;\prod_{l=1}^n\dd\big((\dd/\dd x)^{d_l}\xi^{i_l}\big)(p_l),
\end{equation}
where $\alpha=R.\TFT^{\mathrm{TR}}$ and $\TFT^{\mathrm{TR}}$ is the TFT with $\TFT^{\mathrm{TR}}_{h,n}(e_j,\dots,e_j)=\Delta_j^{(2h-2+n)/2}$ and zero on mixed arguments (the vertex weight of \cite[\S2.3]{DuninBarkowski:2012bw}). We apply it to $f:=4\Wv$, the superpotential \cite[(8-1)]{Dunin-Barkowski:2015caa} with $n=4$ and $a=(0,0,0,0,4v)$: the critical points are $p^5=v$, with $f''(p_c)=20p_c^2\neq 0$, the critical values $5p_c^4$ are pairwise distinct, $\dd y=\dd p\neq 0$, and the shift $p_{\mathrm{Dub}}=p$ is trivial since $a_1=0$; the compatibility condition of \cite[Prop.~4.3 (1)]{Dunin-Barkowski:2015caa} holds, since the curve is rational and $f$ has a pole at $p=\infty$ (moreover $\dd y/\dd f=p^2/(4(p^5-v))$ has poles only at the zeros of $\dd f$); and by the sentence preceding \cite[Thm.~8.1]{Dunin-Barkowski:2015caa} this curve satisfies (2-30)--(2-33), so that \eqref{eq:mod-235} holds for it. The hypotheses on $R$ are \cref{prop:mod-Rtr} at $\kappa=4$.

\emph{Step 2 ($\kappa=4$).} By \cref{prop:mod-Rtr}, for $x=4\Wv$: $R^{4}(\psi)=\Mo(\psi/(4\varsigma^4))$, $\Delta^{4}_j=-20\varsigma_j^2$, $\xi^{(4,j)}=-1/(\sqrt{-20}\,\varsigma_j(p-\varsigma_j))$. So \eqref{eq:mod-235} holds for the correlators $\omega^{(4,1)}_{h,n}$ of $(\bbP^1,4\Wv,p)$ with $\alpha^{4}=R^{4}.\TFT^{4}$.

\emph{Step 3 (transfer to $\kappa=1$).} Take $\sqrt{-20}=2\sqrt{-5}$. Then $w^{4}_j=2w^{1}_j$, $\xi^{(4,j)}=\xi^{(1,j)}/2$, $\dd/\dd x_4=\frac14\dd/\dd x_1$ and $(\Delta^4_j)^{1/2}=2(\Delta^1_j)^{1/2}$; identify the orthonormal bases by $e^4_j\leftrightarrow e^1_j$. Then $\TFT^4_{h,n}=2^{2h-2+n}\TFT^1_{h,n}$, the unit of $\TFT^4$ is half that of $\TFT^1$, and $R^4(\psi)=R^1(\psi/4)$ with $R^1(\psi):=\Mo(\psi/\varsigma^4)$. In the graph sum of $\alpha^4_{h,n}$ a vertex of genus $h_v$ with $n_v$ special points (half-edges, legs, dilaton leaves) gives $2^{2h_v-2+n_v}$; since $\sum_v(2h_v-2+n_v)=2h-2+n+\#\{\text{dilaton leaves}\}$ and each dilaton leaf, linear in the unit, gives $2^{-1}$, the TFT contributes $2^{2h-2+n}$. The substitution $R\mapsto R(\cdot/c)$ multiplies a leg term $\psi^k$ by $c^{-k}$, a dilaton term $r_k\psi^{k+1}$ (of degree $k$ on $\overline{\mathcal M}_{h,n}$ after the leaf is forgotten) by $c^{-k}$, and an edge term $\psi'^a\psi''^b$ (of degree $a+b+1$ on $\overline{\mathcal M}_{h,n}$ after gluing; $E_{R(\cdot/c)}(\psi',\psi'')=c^{-1}E_R(\psi'/c,\psi''/c)$) by $c^{-(a+b+1)}$; hence
\[
[R(\cdot/c).\TFT]_{h,n}=\sum_Dc^{-D}\,[R.\TFT]^{(D)}_{h,n},
\]
where $^{(D)}$ is the part of cohomological degree $D$. So $\alpha^4_{h,n}=2^{2h-2+n}\sum_D4^{-D}[\alpha^1_{h,n}]^{(D)}$ with $\alpha^1:=R^1.\TFT^1$. In \eqref{eq:mod-235} only $D=3h-3+n-\sum_ld_l$ contributes, and the leg $l$ gives $4^{-d_l}2^{-1}$ times the corresponding leg for $\kappa=1$. Hence $\omega^{(4,1)}_{h,n}=2^{2h-2+n}4^{-(3h-3+n)}2^{-n}\,\mathrm{RHS}^1=4^{2-2h-n}\,\mathrm{RHS}^1$, where $\mathrm{RHS}^1$ is the right side of \eqref{eq:mod-235} with the $\kappa=1$ data. By \cref{lem:mod-scale}~(i), $\omega^{(4,1)}_{h,n}=4^{2-2h-n}\omega_{h,n}$. So \eqref{eq:mod-235} holds on $x=\Wv$ with $\alpha^1=\Mo(\cdot/\varsigma^4).\TFT^1$, $\Delta^1_j=-5\varsigma_j^2$ and $\xi^{(1,j)}=-1/(\sqrt{-5}\,\varsigma_j(p-\varsigma_j))$.

\emph{Step 4 (the frozen basis and the factor $(-5)^{h-1}$).} Put $\bar e_j:=e^1_j/\sqrt{-5}$ and $\tilde\xi^j:=\sqrt{-5}\,\xi^{(1,j)}=-1/(\varsigma_j(p-\varsigma_j))$; then $\alpha^1(e^1_{i_1},\dots)\prod\dd(D_x^{d}\xi^{(1,i)})=\alpha^1(\bar e_{i_1},\dots)\prod\dd(D_x^{d}\tilde\xi^{i})$. On the frozen basis $\TFT^1_{h,n}(\bar e_j,\dots,\bar e_j)=(-5)^{-n/2}(\sqrt{-5}\varsigma_j)^{2h-2+n}=(-5)^{h-1}\varsigma_j^{2h-2+n}$, the metric is $\eta(\bar e_j,\bar e_k)=-\delta_{jk}/5$, and the unit is $\hat 1=\sum_k\bar e_k/\varsigma_k$ (\cref{prop:mod-Rtr}~(ii)). Identify $\bar e_j$ with the normalised canonical basis vector $\tilde e_j$ of the model TFT $\TFT^{\mathrm m}_v$, for which $\hat\Delta_j^{1/2}=\zeta^{4j}\varsigma=\varsigma_j$ (\cref{def:set-modcohft}). Then $\TFT^1_{h,n}=(-5)^{h-1}\TFT^{\mathrm m}_{v,h,n}$, with metric $-\frac15$ times that of $\TFT^{\mathrm m}_v$ and the same unit. For a CohFT $\Omega$ with metric $\eta$ and $\lambda\neq 0$, the classes $\lambda^{h-1}\Omega_{h,n}$ form a CohFT for the metric $\eta/\lambda$ with the same unit, and $R.(\lambda^{h-1}\Omega)=\lambda^{h-1}(R.\Omega)$: in the graph sum each vertex gives $\lambda^{h_v-1}$, each edge gives $\lambda$ through the inverse metric, legs and dilaton leaves are unchanged, and $\sum_v(h_v-1)+\#\{\text{edges}\}=h-1$. With $\lambda=-5$: $\alpha^1=(-5)^{h-1}\big(\Mo(\cdot/\varsigma^4).\TFT^{\mathrm m}_v\big)$ on the frozen basis.

\emph{Step 5 (the flat basis).} By \cref{def:set-modcohft}, $\Omega^{\mathrm{mod}}_v=(\Psi^{\mathrm m}_v)_*(\Mo(z/\varsigma^4).\TFT^{\mathrm m}_v)$, and the flat components of the normalised canonical basis are the entries of $\Psi^{\mathrm m}_v$, $\eta(H^b,\tilde e_i)=\varepsilon_b\varsigma^{\nu_b-1}\zeta^{i(b+1)}$ (\cref{def:set-M0}, \cref{prop:set-modrows}~(b): the model row $b$ is $\varsigma^{\nu_b-1}\Rhat^{\mathrm m}{}_b(z/\varsigma^4)$ up to the exponential factor, and $\Rhat^{\mathrm m}(0)=\Psi_0$). Since $H^{3-a}/5$ is the vector $\eta$-dual to $H^a$, $\sum_i\tilde e_i\otimes\tilde\xi^i=\sum_aH^a\otimes\xi'_a$ with $\xi'_a:=\frac15\sum_i\eta(H^{3-a},\tilde e_i)\tilde\xi^i$. Using $\nu_{3-a}=2-\nu_a$, $1/\varsigma_i=\zeta^i/\varsigma$ and $\zeta^{-ia}=(\varsigma_i/\varsigma)^a$,
\[
\xi'_a=-\frac{\varepsilon_{3-a}}{5}\,\varsigma^{-\nu_a-a}\sum_i\frac{\varsigma_i^a}{p-\varsigma_i},\qquad\sum_i\frac{\varsigma_i^a}{p-\varsigma_i}=\begin{cases}5p^4/(p^5-v),&a=0,\\ 5vp^{a-1}/(p^5-v),&1\leq a\leq 4,\end{cases}
\]
by partial fractions ($p^m/(p^5-v)=\sum_i\varsigma_i^{m-4}/(5(p-\varsigma_i))$ for $0\leq m\leq 4$, and $\varsigma_i^5=v$). With $\nu_a+a=0,0,5,5,5$ and $\varepsilon_{3-a}=1,-1,-1,1,1$ for $a=0,\dots,4$ this gives $\xi'_a=\xi_a$ of \cref{def:mod-xi}. Applying $D_x^d$ and $\dd$, \eqref{eq:mod-235} for $x=\Wv$ becomes \eqref{eq:mod-tr}. The sum is finite since a correlator vanishes unless $\sum_ld_l\leq 3h-3+n$. Pairing each slot with a leg, by residues or along rays as in \cref{def:mod-pair} (an identity of forms in each slot), gives the second form. We use neither the identification of $\alpha$ with the CohFT of a Frobenius manifold (the second half of \cite[Thm.~8.1]{Dunin-Barkowski:2015caa}) nor \cite[Prop.~5.2]{Dunin-Barkowski:2015caa}: $\alpha$ is identified with $\Omega^{\mathrm{mod}}_v$ through \cref{prop:mod-Rtr}.
\end{proof}

% Section 6, part 4.
\begin{lem}[Functions of $x$ pair to zero]
\label{lem:mod-xfun}
Let $2h-2+n>0$, $n\geq 1$, let $p_2,\dots,p_n$ be points of $\bbP^1$ that are not branch points, and let $f$ be holomorphic near the five critical values $x(\varsigma_b)$ (for instance $f\in\bbQ(v)[X]$, $v\neq 0$ fixed). Then for each $b$,
\[
\Res_{p=\varsigma_b}f(x(p))\,\omega_{h,n}(p,p_2,\dots,p_n)=0 .
\]
Consequently every pairing of $\omega_{h,n}$ in which one slot is paired with $f(x)$ by the residue pairing, the other slots being paired by residues at the branch points, on circles or along rays, vanishes; legs $\phi,\phi'\in\cL$ with $\phi-\phi'\in\bbQ(v)[x]$ have the same residue pairings; and (with $f=1$) $\omega_{h,n}$ has zero residues at the branch points in each variable.
\end{lem}

\begin{proof}
\emph{The trace identity.} For $X\in\bbP^1$ let $p^1(X),\dots,p^5(X)$ be the points of $x^{-1}(X)$, the roots of $T^5-4XT+4v$, and let $\mathrm{Tr}$ be the push-forward of $1$-forms along $x\colon\bbP^1\to\bbP^1$. We claim that $\mathrm{Tr}\,\omega(X):=\sum_{i=1}^5\omega_{h,n}(p^i(X),p_2,\dots,p_n)=0$ identically. (i) For $X$ not a critical value and $q'\notin x^{-1}(X)$, $\prod_i(q'-p^i(X))=q'^5-4Xq'+4v=4q'\,(x(q')-X)$; since $\omega_{0,2}(p,q')=\dd_p\dd_{q'}\log(p-q')$, this gives $\sum_i\omega_{0,2}(p^i(X),q')=\dd_X\dd_{q'}\log(x(q')-X)=\dd X\,\dd x(q')/(X-x(q'))^2$. (ii) Choose small discs $U_b$ around the branch points with $\iota_b(U_b)=U_b$, disjoint from each other and from $\{p_2,\dots,p_n\}$, and put $U:=\bigcup_bU_b$. For $p_1\notin U$ the recursion \eqref{eq:mod-kernel} reads $\omega_{h,n}(p_1,J)=\sum_b(2\pi\ri)^{-1}\oint_{\partial U_b}K_b(p_1,q)\,\mathcal R_b(q,J)$, $J=(p_2,\dots,p_n)$, where $\mathcal R_b$ (a local letter) is the bracket of \eqref{eq:mod-kernel}, the factors $\omega_{0,2}(q,p_l)$ included; it does not depend on $p_1$, and the integrand has no pole in $U_b$ other than $\varsigma_b$. The path from $\iota_bq$ to $q$ in $K_b$ stays in $U_b$. (iii) For $X$ outside the compact set $x(\overline U)$ all $p^i(X)$ lie outside $U$, and by (i)
\[
\sum_i\int_{\iota_bq}^q\omega_{0,2}(p^i(X),\cdot)=\dd X\Big[\frac{1}{X-x(q)}-\frac{1}{X-x(\iota_bq)}\Big]=0,
\]
because $x(\iota_bq)=x(q)$. Summing (ii) over $i$ (a finite sum, exchanged with the contour integrals) gives $\mathrm{Tr}\,\omega(X)=0$ for these $X$; as $\mathrm{Tr}\,\omega$ is a meromorphic $1$-form in $X$, it vanishes identically. (This is the computation of \cite[App.~A, (A-1), (A-5)]{Eynard:2007kz}.)

\emph{The residues.} Fix $b$ and put $X_b:=x(\varsigma_b)$. The fibre $x^{-1}(X_b)$ consists of $\varsigma_b$, of local degree $2$, and three further points $q_1,q_2,q_3$ of local degree $1$; none of them is a branch point, because the critical values are pairwise distinct, and none is $0$ or $\infty$, which are mapped to $\infty$. Hence $f(x)\,\omega_{h,n}(\cdot,J)$ is holomorphic at the $q_k$ (\cref{conv:mod-kernel}: poles only at branch points, none on the diagonals). For a $1$-form $\eta$ meromorphic near the fibre, $\Res_{X=X_b}\mathrm{Tr}(\eta)=\sum_{q\in x^{-1}(X_b)}\Res_q\eta$: the preimage of a small circle around $X_b$ is a union of small loops, one around each point of the fibre, mapped onto the circle with the local degree, so $\oint\mathrm{Tr}(\eta)=\sum_q\oint\eta$. With $\eta=f(x)\,\omega_{h,n}(\cdot,J)$ we have $\mathrm{Tr}(\eta)=f(X)\,\mathrm{Tr}\,\omega=0$, so $0=\Res_{\varsigma_b}\eta+\sum_k\Res_{q_k}\eta=\Res_{\varsigma_b}\eta$.

\emph{Consequences.} $\omega_{h,n}$ is symmetric (the paragraph after \cref{conv:mod-kernel}), so the slot carrying $f(x)$ may be taken to be the first. The contours of the other slots can be chosen disjoint from small discs around the branch points used in the first slot, and $\omega_{h,n}$ has no diagonal poles, so the residue in the first slot can be taken first, for fixed points of the other contours, and it vanishes. The other two statements follow.
\end{proof}

The statement and its deduction are also \cite[Thm.~4.4 (4-23), Cor.~4.1 (4-26)--(4-29), App.~A (A-25)]{Eynard:2007kz}; the proof above does not use them.

\subsection{The leg dictionary and the residue identity}
\label{sec:mod-legs}

\begin{defn}[The legs]
\label{def:mod-legs}
Let $\phi^b_0:=p,\ \log p,\ v/p,\ p^3,\ p^2$ for $b=0,\dots,4$. For $F$ in the span of the $p^j(\log p)^l$ ($j\in\bbZ$, $l\geq 0$) with coefficients in $\bbQ[v]$, let $I(F)$ be the termwise antiderivative of $-F\,\dd x=-F(p^3-vp^{-2})\,\dd p$ with integration constant $0$: the termwise antiderivative of $p^j(\log p)^l\dd p$ is $p^{j+1}$ times a polynomial in $\log p$ if $j\neq -1$, and $(\log p)^{l+1}/(l+1)$ if $j=-1$, so it has no term $p^0(\log p)^0$, and $I(F)$ is the antiderivative without such a term. Put $\phi^b_{k+1}:=I(\phi^b_k)$. Then $D_x\phi^b_{k+1}=-\phi^b_k$, and $\phi^b_k$ is homogeneous of weight $\wt(\phi^b_k)=w_b+4k$, where $w_b:=\nu_b+1$, that is $(w_0,\dots,w_4)=(1,0,4,3,2)$.
\end{defn}

\begin{prop}[The leg dictionary]
\label{prop:mod-LD}
Let $\varrho=(\varrho_0,\dots,\varrho_4):=(-1,1,-1,-3,-2)$. For all $a,b\in\bbZ/5$ and $k,d\geq 0$,
\begin{equation}
\label{eq:mod-LD}
\langle\phi^b_k,\dd\xi_{a,d}\rangle=\sum_c\Res_{c}\phi^b_k\,\dd\xi_{a,d}=-\varrho_b\,\delta_{ab}\,\delta_{dk},
\end{equation}
$c$ running over the branch points and $\log p$ being any branch near each of them. Hence $q_{\phi^b_k}=-\varrho_b\,z^kH^b$, and the TR leg of $v^mz^kH^b$ (the insertion $v^m\psi^kH^b$, \cref{conv:tra-forms}) is $-\varrho_b^{-1}v^m\phi^b_k$. In particular the unit $H^0$ has the TR leg $y=p$; $\Phi_{y\dd x}=p^5/5-v\log p=-\phi^0_1$ has $q_{\Phi_{y\dd x}}=-z\,H^0$; $-\log p$ has $q_{-\log p}=H$; the Euler-dilaton vector $\se=z\,H^0+vH$ has the TR leg $-p^5/5$; and $q_{f(x)}=0$ for $f\in\bbQ(v)[x]$.
\end{prop}

\begin{proof}
\emph{Step 1 (branches and functions of $x$).} Each leg is a polynomial in $\log p$ with coefficients in $\bbQ[v][p,p^{-1}]$, so it is analytic near every branch point $c\neq 0$, for any branch of $\log p$. A change of branch $\log p\mapsto\log p+2\pi\ri N$ changes $\phi^b_k$ by a function $\delta^b_k$ (a local letter) that is a polynomial in $x$ with constant coefficients: $\delta^b_0$ is $0$ or $2\pi\ri N$; the identity $\dd\phi^b_{k+1}=-\phi^b_k\,\dd x$ of \cref{def:mod-legs} is an identity of expressions in $p$ and $\log p$, with $\dd\log p=\dd p/p$, and it survives the substitution, so $\dd\delta^b_{k+1}=-\delta^b_k\,\dd x$; and if $\delta^b_k=P(x)$ with a polynomial $P$, then $\delta^b_{k+1}+\int_0^xP$ has zero differential near $c$, hence is constant. For $f$ holomorphic near $c$, repeated integration by parts gives $\Res_cf(x)\,\dd\xi_{a,d}=\pm\Res_cf^{(d+1)}(x)\,\xi_a\dd x$, which vanishes because $\xi_a\dd x=\varepsilon'_av^{\alpha_a}p^{m_a-2}\dd p$ is holomorphic at $c\neq 0$. So the pairings do not depend on the branches, and $q_{f(x)}=0$.
\emph{Step 2 (lowering).} For $d\geq 1$, integration by parts at each $c$ gives $\langle\phi,\dd\xi_{a,d}\rangle=-\sum_c\Res_c\xi_{a,d}\,\dd\phi=-\langle D_x\phi,\dd\xi_{a,d-1}\rangle$, since $\xi_{a,d}\dd\phi=(D_x\xi_{a,d-1})(D_x\phi)\,\dd x=(D_x\phi)\,\dd\xi_{a,d-1}$. Hence $\langle\phi^b_{k+1},\dd\xi_{a,d}\rangle=\langle\phi^b_k,\dd\xi_{a,d-1}\rangle$.
\emph{Step 3 ($d=0$, $k\geq 1$).} $\langle\phi^b_k,\dd\xi_a\rangle=-\sum_c\Res_c\xi_a\,\dd\phi^b_k=\sum_c\Res_c\phi^b_{k-1}\,\xi_a\dd x=0$, since $\xi_a\dd x$ is holomorphic at $c$.
\emph{Step 4 ($k=0$, $d\geq 1$).} $\langle\phi^b_0,\dd\xi_{a,d}\rangle=-\sum_c\Res_cD_x^d(\xi_a)\,\dd\phi^b_0$. This form is meromorphic on $\bbP^1$ ($\dd\log p=\dd p/p$); it is holomorphic at $p=0$, because $D_x$ maps functions holomorphic at $0$ to $O(p^2)$ while $\dd\phi^b_0=O(p^{-2})\dd p$; and it is $O(p^{1-4d})\dd p$ at $\infty$, because $D_x^d\xi_a=O(p^{-1-4d})$ and $\dd\phi^b_0=O(p^2)\dd p$. By the residue theorem the residues at the branch points add up to $0$.
\emph{Step 5 ($k=d=0$).} Write $\dd\phi^b_0=\gamma_bp^{n_b}\dd p$ with the local constants $(\gamma_b,n_b)=(1,0),(1,-1),(-v,-2),(3,2),(2,1)$. Since $\Res_cp^N\dd p/(p^5-v)=c^{N-4}/5$ and $\sum_cc^N=5v^{N/5}$ if $5\mid N$ and $0$ otherwise,
\[
\langle\phi^b_0,\dd\xi_a\rangle=-\varepsilon'_a\gamma_bv^{\alpha_a}\,\frac15\sum_cc^{m_a+n_b-4},
\]
which is non-zero only if $m_a+n_b\equiv 4\pmod 5$, that is, only if $a=b$. For $a=b=0,\dots,4$ the values are $1,-1,1,3,2$, that is $-\varrho_b$ (for $b=1$: $-v\cdot v^{-1}$; for $b=2$: $-(-v)v^{-1}$).
\emph{Step 6.} If $k=d$, Step 2 applied $k$ times and Step 5 give $-\varrho_b\delta_{ab}$; if $k>d$, Step 2 and Step 3 give $0$; if $k<d$, Step 2 and Step 4 give $0$. This is \eqref{eq:mod-LD}, and $q_{\phi^b_k}=\sum_{a,d}\langle\phi^b_k,\dd\xi_{a,d}\rangle H^az^d=-\varrho_bz^kH^b$. The consequences: $\phi^0_0=p$ and $-\varrho_0=1$; $\phi^0_1=I(p)=-(p^5/5-v\log p)$, so $q_{\Phi_{y\dd x}}=-q_{\phi^0_1}=\varrho_0z\,H^0=-zH^0$; $q_{\log p}=-\varrho_1H=-H$; and the TR leg of $\se$ is $\phi^0_1-v\phi^1_0=-p^5/5$.
\end{proof}

\begin{cor}[Free energies of the model and of the curve]
\label{cor:mod-trF}
For $h\geq 2$ let $\Fm_h(v)=\langle\ \rangle_h(v)=\int_{\overline{\mathcal M}_h}\Omega^{\mathrm{mod}}_{v;h,0}$ be the genus-$h$ free energy of the model CohFT (\cref{def:set-modcohft}), the value at the base point $-z1$ of the dilaton-shifted potential $\overline F_h$ (\cref{def:set-modcohft,def:set-graphsum}). Then $\Fm_h=(-5)^{1-h}\FTR_h$.
\end{cor}

\begin{proof}
On the side of the curve, $(2-2h)\FTR_h=\sum_b\Res_{\varsigma_b}\Phi_{y\dd x}\,\omega_{h,1}=\langle\Phi_{y\dd x}\rangle^{\mathrm{TR}}_h$ by \cref{conv:mod-kernel}; the branch of $\log p$ does not matter, since $\omega_{h,1}$ has zero residues (\cref{lem:mod-xfun}). Only this definition is used on the side of the curve. (It is the definition \cite[Def.~4.3 (4-10)]{Eynard:2007kz} up to the sign $\FTR_h=-F^{(h)}$, and, as remarked after \cite[Thm.~4.7]{Eynard:2007kz}, it extends to $n=0$ the dilaton equation \cite[Thm.~4.7 (4-30)]{Eynard:2007kz}, stated there for $n\geq 1$, which in the convention of \cref{conv:mod-kernel} reads $\sum_b\Res_{\varsigma_b}\Phi_{y\dd x}\,\omega_{h,n+1}(\cdot,J)=(2-2h-n)\,\omega_{h,n}(J)$ for $2h-2+n>0$.) By \cref{thm:mod-tr} and \cref{prop:mod-LD}, $\langle\Phi_{y\dd x}\rangle^{\mathrm{TR}}_h=(-5)^{h-1}\langle q_{\Phi_{y\dd x}}\rangle_h=-(-5)^{h-1}\langle\psi H^0\rangle_h$. The model CohFT has the flat unit $H^0$, and the dilaton equation (\cref{lem:set-dilaton}~(i), (ii), with no further insertion) gives $\langle\psi H^0\rangle_h=(2h-2)\Fm_h$. Comparing, $(2-2h)\FTR_h=(2-2h)(-5)^{h-1}\Fm_h$.
\end{proof}

% Section 6, part 5.
The weight spaces and the modules of legs are needed before the logarithmic part of the legs can be computed; they are used again in \cref{sec:mod-lattice}.

\begin{defn}[Weight spaces, the leg map, the leg class]
\label{def:mod-weights}
For $w\geq 0$ let $V_w$ be the $\bbQ$-span of the insertions $v^m\psi^kH^b$ ($m,k\geq 0$, $b\in\bbZ/5$) with $w_b+4k+5m=w$, and $\mathbb F:=\bigoplus_wV_w$, the free graded $\bbQ[v]$-module on the $\psi^kH^b$ (the weights of \cref{conv:tra-forms}). The \emph{leg map} is the $\bbQ$-linear map $\Lambda\colon\mathbb F\to\{\text{legs}\}$, $\Lambda(v^m\psi^kH^b):=\varrho_b^{-1}v^m\phi^b_k$; by \cref{prop:mod-LD} the TR leg of $X\in\mathbb F$ is $-\Lambda(X)$. Let $\bbQ[p,\bar p]$, $\bar p=v/p$, be the subring of $\bbQ[v][p,p^{-1}]$ spanned by the $v^mp^j$ with $j\geq -m$; recall the leg class $\cL=\bbQ[v][p]+\bbQ(v)[x]$ (\cref{def:mod-pair}). Put
\[
\mathsf M:=\bbQ[p,\bar p]\oplus\log p\cdot\bbQ[v,x],\qquad \mathsf N:=\mathsf M/\bbQ[v,x]=\bbQ[p,\bar p]/\bbQ[v,x]\oplus\log p\cdot\bbQ[v,x],
\]
(the sum is direct since $1$ and $\log p$ are independent over Laurent polynomials; $\bbQ[v,x]\subset\bbQ[p,\bar p]$ as $v=p\bar p$, $x=p^4/4+\bar p$), graded by weight. For $F=A+\log p\cdot B\in\mathsf M$ the element $B\in\bbQ[v,x]$ is the \emph{$\log p$-coefficient} of $F$, and $F$ is \emph{log-free} if $B=0$. Finally $\mathrm d(w):=\lfloor w/5\rfloor+1-[5\mid w]$ and $\mathrm e(w):=\#\{(m,n)\in\bbZ_{\geq 0}^2: 5m+4n=w\}$.
\end{defn}

\begin{lem}[Closure]
\label{lem:mod-closure}
\begin{enumerate}
\item For $a,c\geq 0$, $I(p^a\bar p^c)=-p^{a+4}\bar p^c/(a-c+4)+p^a\bar p^{c+1}/(a-c-1)$, a term with a zero denominator being read as $-v^c\log p$, respectively $+v^{c+1}\log p$.
\item For $g\in\bbQ[v,x]$ with an $x$-primitive $G$, $I(\log p\cdot g(x))=-\log p\cdot G(x)+\int G(x)\,\dd p/p$, where the last integral is termwise: $\int p^a\bar p^c\,\dd p/p=p^a\bar p^c/(a-c)$ if $a\neq c$ and $=v^c\log p$ if $a=c$.
\item $I$ is $\bbQ[v]$-linear, $I(\mathsf M)\subset\mathsf M$ and $I(\bbQ[v,x])\subset\bbQ[v,x]$. Hence every leg $\phi^b_k$ lies in $\mathsf M$, and $\mathsf M$ is a $\bbQ[v,x]$-module.
\end{enumerate}
\end{lem}

\begin{proof}
(i) $-\int v^cp^{a-c}(p^3-vp^{-2})\,\dd p=-v^cp^{a-c+4}/(a-c+4)+v^{c+1}p^{a-c-1}/(a-c-1)$, with $\int\dd p/p=\log p$ in the exceptional cases; and $v^cp^{a-c+4}=p^{a+4}\bar p^c$, $v^{c+1}p^{a-c-1}=p^a\bar p^{c+1}$. (ii) $\dd(\log p\cdot G(x))=G(x)\,\dd p/p+\log p\cdot g(x)\,\dd x$, and $G(x(p))\in\bbQ[p,\bar p]$; the right side has no term $p^0(\log p)^0$ (the termwise $\int G\,\dd p/p$ has none, and $\log p\cdot G$ has no term free of $\log p$), so it is the antiderivative of \cref{def:mod-legs}. A different primitive $G+C$ changes $-\log p\cdot G$ by $-C\log p$ and $\int G\,\dd p/p$ by $C\log p$, so the identity holds for every primitive. (iii) Linearity is clear. By (i), $I$ maps $\bbQ[p,\bar p]$ into $\mathsf M$, the exceptional terms lying in $\log p\cdot\bbQ[v]$. By (ii), $I(\log p\cdot g)\in\mathsf M$, since $G\in\bbQ[v,x]$ and $\int G\,\dd p/p\in\bbQ[p,\bar p]+\log p\cdot\bbQ[v]$. For $g\in\bbQ[v,x]$, $-g(x)\,\dd x=-\dd G(x(p))$ is exact, with no term $p^{-1}\dd p$, so $I(g)=-G(x)+[p^0]G(x(p))\in\bbQ[v,x]$. The five $\phi^b_0$ lie in $\mathsf M$, and induction gives $\phi^b_k\in\mathsf M$. Finally $\bbQ[p,\bar p]$ is a ring containing $\bbQ[v,x]$.
\end{proof}

\begin{lem}[Generating function of the log coefficients]
\label{lem:mod-genfn}
For all $b$ and $k$, the $\log p$-coefficient of $\phi^b_k$ is $\sum_{i=0}^k\mathcal R^b_{k-i}(-x)^i/i!$, where $\mathcal R^b_n:=\Res_{p=0}\Wv^n\,\dd\phi^b_0/n!$, a monomial in $v$.
\end{lem}

\begin{proof}
Fix $b$, write $\phi_k:=\phi^b_k$, and put $\mathcal B_n:=\sum_{j=0}^nx^j\phi_{n-j}/j!$, the coefficient of $\epsilon^n$ in $e^{\epsilon x}\sum_k\epsilon^k\phi_k$ ($\epsilon$ a formal variable). Since $\dd\phi_m/\dd p=-x'\phi_{m-1}$ for $m\geq 1$, the derivative telescopes: $\dd\mathcal B_n/\dd p=x^n\phi_0'/n!$, a Laurent polynomial in $p$. By \cref{lem:mod-closure}~(iii), $\mathcal B_n\in\mathsf M$; write $\mathcal B_n=A+\log p\cdot B$ with $A\in\bbQ[p,\bar p]$ and $B\in\bbQ[v,x]$. Then $\dd\mathcal B_n/\dd p=A'+B/p+\log p\cdot\dd B(x(p))/\dd p$ has no logarithmic term, so $B(x(p))$ does not depend on $p$, and $B\in\bbQ[v]$. As $A'$, a derivative, has no term $p^{-1}$, $B$ is the coefficient of $p^{-1}$ in $x^n\phi_0'/n!$, that is $B=\mathcal R^b_n$. Inverting, $\sum_k\epsilon^k\phi_k=e^{-\epsilon x}\sum_n\epsilon^n\mathcal B_n$, so $\phi_k=\sum_{i=0}^k(-x)^i\mathcal B_{k-i}/i!$; taking $\log p$-coefficients, which commutes with multiplication by $\bbQ[v,x]$, gives the claim. By weight, $\mathcal R^b_n$ is a monomial in $v$.
\end{proof}

\begin{lem}[The residue identity]
\label{lem:mod-resid}
For all $n\geq 0$ and $b$, $\mathcal R^b_n=\varrho_b\,\Pi_0(\psi^nH^b)$.
\end{lem}

\begin{proof}
By \cref{def:gau-van}, $\IC_m=\Res_{p=0}p^me^{-\Wv(p)/z}\dd p$ is a formal residue, taken coefficientwise, so
\[
[z^{-n}]\IC_m=\frac{(-1)^n}{n!}\Res_{p=0}p^m\Wv^n\,\dd p .
\]
With $\YC_b=\varepsilon_b\IC_{\nu_b}$ and $\Pi_0(z^nH^b)=(-1)^{n+1}[z^{-n}]\YC_b$ (\cref{def:tra-tower}) this gives
\[
\Pi_0(\psi^nH^b)=-\frac{\varepsilon_b}{n!}\Res_{p=0}p^{\nu_b}\Wv^n\,\dd p .
\]
Compare with $\dd\phi^b_0=\dd p,\ \dd p/p,\ -vp^{-2}\dd p,\ 3p^2\dd p,\ 2p\,\dd p$: for $b=0,1,3,4$ the identity is immediate, with $\varrho_b=-1,1,-3,-2$. For $b=2$, $\Res_{p=0}\Wv^n(p^3-vp^{-2})\,\dd p=\Res_{p=0}\dd(\Wv^{n+1})/(n+1)=0$, so $\mathcal R^2_n=-v\Res_{p=0}p^{-2}\Wv^n\dd p/n!=-\Res_{p=0}p^3\Wv^n\dd p/n!=-\Pi_0(\psi^nH^2)$, and $\varrho_2=-1$.
\end{proof}

\begin{prop}[Log coefficients of the legs]
\label{prop:mod-dii}
For all $b$ and $k$,
\begin{equation}
\label{eq:mod-dii}
\text{the }\log p\text{-coefficient of }\phi^b_k\ =\ \varrho_b\sum_{i=0}^k\Pi_i(\psi^kH^b)\,\frac{(-x)^i}{i!}.
\end{equation}
Hence, for $X\in V_w$, the $\log p$-coefficient of $\Lambda(X)$ is $\sum_{i\geq 0}\Pi_i(X)(-x)^i/i!$, and $\Lambda(X)$ is log-free if and only if $X\in\sKo$.
\end{prop}

\begin{proof}
By \cref{lem:mod-genfn,lem:mod-resid} and the shift rule $\Pi_i(\psi^kH^b)=\Pi_0(\psi^{k-i}H^b)$ of \cref{def:tra-tower}, the $\log p$-coefficient of $\phi^b_k$ is $\sum_i\varrho_b\Pi_0(\psi^{k-i}H^b)(-x)^i/i!=\varrho_b\sum_i\Pi_i(\psi^kH^b)(-x)^i/i!$. For $X=\sum_ta_tv^{m_t}\psi^{k_t}H^{b_t}\in V_w$ the factors $\varrho_b^{-1}$ of $\Lambda$ cancel, and since the $\Pi_i$ are $\bbQ[v]$-linear the $\log p$-coefficient of $\Lambda(X)$ is $\sum_i\Pi_i(X)(-x)^i/i!$. The $\Pi_i(X)$ lie in $\bbQ[v]$ and the powers of $x$ are linearly independent over $\bbQ(v)$, so $\Lambda(X)$ is log-free if and only if $\Pi_i(X)=0$ for all $i\geq 0$, that is (\cref{def:tra-slice}) if and only if $X\in\sKo$.
\end{proof}

The normalising constants $\varrho_b$ in $\Lambda$ matter for $\sKo$ as a subspace: $\se=z\,H^0+vH$ lies in $\sKo$ and $\Lambda(\se)=p^5/5$ is log-free, whereas $\phi^0_1+v\phi^1_0=-p^5/5+2v\log p$ is not.

\input{sections/s06_model_6}



